\section{Rational matrix quotients without simultaneous diagonalization}
\label{sec:general-quotient}
The normalization operator is not peculiar to the stochastic diagonal
model. We state its exact relation to scalar-denominator rational
interpolation, and then separate the general resolvent bound from the
stronger additive law furnished by the joint coefficient algebra.

\begin{theorem}[The rational interpolation module]\label{thm:rational-module}
Let $K$ be any rectangular matrix polynomial of degree at most $d$,
and let $q$ be monic of degree $d$, with $q(T_j)\ne0$. Form $\cN_P$
from $P(T_j)=K(T_j)/q(T_j)$ at $\ell\ge2d+1$ distinct nodes. Then
\begin{equation}\label{eq:general-gcd}
 \ker\cN_P=\{(q/h)s:\deg s<\deg h\},\qquad
 h=\gcd(q,\hbox{all entries of }K).
\end{equation}
More generally, the pairs $(H,r)$ with $\deg H\le d$, $\deg r<d$ and
$H(T_j)=r(T_j)P(T_j)$ are exactly
$((K/h)s,(q/h)s)$ with $\deg s<\deg h$.
No diagonalization, channel factorization or positivity is assumed.
\end{theorem}
\begin{proof}
Interpolation gives $qH-rK=0$ at $2d+1$ nodes. Each entry has degree
at most $2d$, so this is an identity. The polynomial $q/h$ and the
entries of $K/h$ have no common divisor. A Bezout identity for this
collection, or prime-by-prime divisibility, shows that $q/h$ divides
$r$. Write $r=(q/h)s$ and then $H=(K/h)s$. The degree constraint on
$r$ is exactly $\deg s<\deg h$. Conversely every such pair satisfies
the identity and the degree constraints.
\end{proof}

In the probability model $q$ is the sum of the entries of $K$, so the
common divisor of those entries already divides $q$. Full-rank constant
channels then identify it with $\gcd(f_b)$, giving
Lemma~\ref{lem:defect}. The module itself is an instance of classical
simultaneous rational interpolation, not a new interpolation algorithm.
For any monic $h_u=h+us$ of the same degree, the pair
$((K/h)h_u,(q/h)h_u)$ has exactly the same rational curve. In a square
regular numerator the cancelled scalar factor contributes its zeros
with multiplicity equal to the matrix size. This identifies the
spectral information absent from a minimal realization.

Here is a local observable spectral theorem for general square
numerators. It genuinely removes simultaneous diagonalization, though
it does not impose the diagonal model's better normalization constant.
Let $M=K$ have invertible leading coefficient, let $\tau(P)>0$, and let
$X$ be its finite root set. On a fixed coefficient-bounded neighbourhood
assume the clock denominators are bounded away from zero. Let $m$ be
the largest pole order of $M^{-1}$. Write
\begin{align*}
 M(z)^{-1}&=\sum_{x\in X}\sum_{j=1}^m\frac{R_{x,j}}{(z-x)^j},\\
 \mathcal T(z)r:=M(z)^{-1}I_d(rP)(z)
   &=T_\infty r+\sum_{x\in X}\sum_{j=1}^m\frac{T_{x,j}r}{(z-x)^j},\\
 b&=\|M_d^{-1}\|+\sum_{x,j}\|R_{x,j}\|,\qquad
 t=\|T_\infty\|+\sum_{x,j}\|T_{x,j}\|.
\end{align*}
The $T$ coefficients are linear maps from the Euclidean coefficient
space of $r$ to matrices with operator norm. All coefficients are
observable Laurent coefficients of the recovered rational matrix.

\begin{proposition}[General quotient resolvent bound]\label{prop:general-bound}
For sufficiently close degree-$d$ competitors with regular leading
coefficient, their $kd$ numerator roots satisfy
\[
 d_\infty(Z(M),Z(M'))\le C\{\delta(b+t/\tau(P))\}^{1/m}.
\]
Complex roots and nondiagonalizable matrix coefficients are allowed.
\end{proposition}
\begin{proof}
Write $r=q'-q$. The normalization identity gives $\|r\|\le C\delta/\tau$,
and $M'=M+I_d(rP)+H$ with $\|H\|_{\rm coef}\le C\delta$.
For $\rho=\dist(z,X)\le1$ on a fixed disk, the displayed Laurent
expansions imply
\[
 \|M(z)^{-1}(I_d(rP)+H)(z)\|
 \le C\rho^{-m}(t\|r\|+b\delta).
\]
The homotopy is nonzero outside disks of radius
$C(t\|r\|+b\delta)^{1/m}$. Coefficient closeness preserves its leading
invertibility and excludes roots outside a fixed disk. Argument-principle
counting on each disk component, followed by the finite-permutation
limit, proves the assertion including algebraic multiplicities.
\end{proof}

\subsection{Comparison with realization and matrix interpolation}
Anderson--Antoulas \cite{AA90} recover minimal state-variable
realizations from rational matrix interpolation data through Loewner
matrices. Beckermann--Labahn \cite{BL00} compute matrix rational
interpolants and matrix greatest common divisors, including singular
interpolation subproblems, by fraction-free structured algorithms.
Mayo--Antoulas \cite{MA07} develop the generalized realization
framework. These are direct predecessors of the normalization module,
not merely neighbouring applications. In particular
Theorem~\ref{thm:rational-module} locates our linear system in their
interpolation setting: it is the bounded-degree, scalar-denominator
submodule obtained after eliminating numerator coefficients by an
orthogonal Vandermonde residual.

The statistical problem asks for something different from the minimal
realization alone. A common factor changes the latent numerator
spectrum while leaving the entire rational datum unchanged. Specifying
the original degree and stochastic realization does not undo that
cancellation. Lemma~\ref{lem:local-fibre} shows that the ambiguity is
realized by positive fixed-channel parameters, not just by formal
rational representations. Proposition~\ref{prop:full-span} distinguishes
the full-complexity clock count from the worst-case scalar obstruction.

For a given polynomial matrix, the exclusion and multiplicity argument
is the classical mechanism of Bauer--Fike and its polynomial and
pencil extensions \cite{BF,Chu87,Chu03,TH,HMT,SB}. We use it explicitly
in Lemma~\ref{lem:roots} and Proposition~\ref{prop:general-bound}.
The further assertion of Theorem~\ref{thm:joint} is the conversion
from normalized probability error to that spectral perturbation: the
joint coefficient projectors keep the normalization term scalar on
each complete-polynomial group, yielding $B_d+\tau^{-1}$ instead of
a second multiplication by channel conditioning. The exact metric
identity in Theorem~\ref{thm:exact-condition} and the positive testing
families below then address the target and the observation experiment,
not the numerical complexity of rational interpolation. No claim that
Loewner or matrix-GCD methods fail on these data is needed or made.
